@ID: ON18D1P1U
@BIDANG: Analisis Kompleks




Diberikan bilangan asli $n$. Cari semua bilangan bulat yang dapat dinyatakan dalam bentuk

$$
\left(\sum_{k=1}^{n} z_k\right)
\left(\sum_{k=1}^{n}\frac1{z_k}\right)
$$

untuk suatu bilangan kompleks $z_1,\dots,z_n\in\mathbb C$ dengan $|z_1|=\cdots=|z_n|=1$.

@ID: ON18D1P2U
@BIDANG: Aljabar Linear

Misalkan $A\in\mathbb R^{n\times n}$ dengan $k_1,k_2,\dots,k_n\in\mathbb R^n$ adalah vektor-vektor kolom dari $A$. Vektor-vektor kolom tersebut memenuhi hubungan

$$
k_i=(i+2)k_{i+2},\qquad i=1,2,\dots,n-2.
$$

Untuk $n>3$, pilihlah satu nilai eigen $A$ kemudian tentukan dimensi terkecil yang mungkin untuk ruang eigen dari nilai eigen yang terpilih.

@ID: ON18D1P3U
@BIDANG: Kombinatorika

Untuk bilangan bulat positif $n\ge1$, definisikan $h_n$ sebagai banyaknya cara menuliskan $n$ sebagai jumlahan dari bilangan $1$ dan $2$ dengan memperhatikan urutan kemunculan bilangan $1$ dan $2$. Sebagai contoh, $h_3=3$, karena

$$
3=1+1+1,\qquad
3=2+1,\qquad
3=1+2.
$$

Selanjutnya, $h_4=5$, karena

$$
4=1+1+1+1,\;
4=1+1+2,\;
4=1+2+1,\;
4=2+1+1,\;
4=2+2.
$$

Buktikan bahwa untuk $k\ge0$

$$
h_{2k+1}
=
\sum_{i\ge0}\sum_{j\ge0}
\binom{k-i}{j}
\binom{k-j}{i}.
$$

@ID: ON18D1P4U
@BIDANG: Analisis Real

Diberikan fungsi kontinu $f:[0,\infty)\to[0,a]$, untuk suatu bilangan real $a>0$, dan $f$ memenuhi *the intermediate value property* pada $[0,\infty)$. Jika $f(0)=0$ dan

$$
xf(x)\ge\int_0^x f(t)\,dt,\qquad \forall x\in(0,\infty),
$$

buktikan bahwa $f$ mempunyai antiderivatif.

@ID: ON18D1P5U
@BIDANG: Struktur Aljabar

Diberikan bilangan prima $p$ dan bilangan asli $n\ge3$. Misalkan

$$
G_n(p)=
\left\{
x+a_2x^2+\cdots+a_nx^n
\mid
a_i\in\mathbb Z_p
\right\}.
$$

Definisikan operasi $\circ$ di $G_n(p)$ melalui

$$
P\circ Q=P(Q(x))
\pmod{x^{n+1}}.
$$

Sebagai contoh di $G_3(5)$, misalkan $P=x+2x^2$ dan $Q=x+3x^2$. Maka

$$
\begin{aligned}
P\circ Q
&=(x+3x^2)+2(x+3x^2)^2\\
&=x+3x^2+2(x^2+6x^3+9x^4)\\
&=x+5x^2+12x^3+18x^4\\
&=x+5x^2+12x^3
\qquad(\text{karena modulo }x^4)\\
&=x+2x^3
\qquad(\text{karena koefisiennya di }\mathbb Z_5).
\end{aligned}
$$

Periksa apakah $G_n(p)$ merupakan grup terhadap operasi $\circ$.

@ID: ON18D2P1U
@BIDANG: Kombinatorika

Diberikan bilangan bulat tak negatif $k$ dan $n$ sehingga $0\le k<n$. Berikan bukti kombinatorial bahwa

$$
\sum_{j=0}^{k}\binom{n}{j}
=
\sum_{j=0}^{k}\binom{n-1-j}{k-j}2^j.
$$

@ID: ON18D2P2U
@BIDANG: Struktur Aljabar

(a) Apakah $\langle x+2018\rangle$ merupakan ideal prima di $\mathbb Z[x]$?

(b) Apakah $\langle x+2018\rangle$ merupakan ideal maksimal di $\mathbb Z[x]$?

@ID: ON18D2P3U
@BIDANG: Analisis Real

Misalkan barisan bilangan real $\{x_n\}$ terbatas dan memenuhi

$$
x_{n+2}\le\frac12(x_n+x_{n+1}),\qquad \forall n\ge0.
$$

Jika $A_n=\max\{x_n,x_{n+1}\}$, buktikan:

(a) barisan $\{A_n\}$ konvergen.

(b) barisan $\{x_n\}$ konvergen.

@ID: ON18D2P4U
@BIDANG: Analisis Real

Buktikan atau beri contoh penyangkal.

(a) Untuk setiap fungsi $f:\mathbb R\to\mathbb R$, jika fungsi $g:\mathbb R\to\mathbb R$ dengan definisi $g(x)=f(x^{2018})$ untuk setiap $x\in\mathbb R$, mempunyai deret pangkat

$$
\sum_{n=0}^{\infty}a_nx^n
$$

yang konvergen di suatu persekitaran $x=0$, maka $f$ mempunyai turunan di $x=0$.

(b) Untuk setiap fungsi $f:\mathbb C\to\mathbb C$, jika fungsi $g:\mathbb C\to\mathbb C$ dengan definisi $g(z)=f(z^{2018})$ untuk setiap $z\in\mathbb C$, mempunyai deret pangkat

$$
\sum_{n=0}^{\infty}a_nz^n
$$

yang konvergen di suatu persekitaran $z=0$, maka $f$ mempunyai turunan di $z=0$.

@ID: ON18D2P5U
@BIDANG: Aljabar Linear

Diberikan bilangan asli $n$. Misalkan

$$
A=
\begin{pmatrix}
a&b\\
c&d
\end{pmatrix}
\in\mathbb R^{2\times2}
$$

dengan $\det(A)=a+d=1$. Jika

$$
A^n=
\begin{pmatrix}
a'&b'\\
c'&d'
\end{pmatrix},
$$

tentukan nilai $a'+d'$.